\section{Related Work}

\subsection{Threat-intelligence exchange}

STIX/TAXII, MISP, and DShield distribute threat information \cite{stix21,taxii21,mispCore,dshieldFeeds}. TCOP may ingest evidence derived from them; it addresses the later agent-runtime question of whether a receiver should authorize one pending tool or capability action after receipt-, session-, and action-specific validation. These systems do not specify that action-time receiver-local authorization contract.

\subsection{Response orchestration}

RID, DOTS, and OpenC2 exchange incident reports, mitigation signals, and defense commands \cite{rid12,dots21,openc2v1}. CACAO and SOAR coordinate playbooks and response workflows \cite{cacao20,bridges23}. Either can implement a response after TCX acceptance. TCOP instead bounds the preceding decision input: whether signed runtime evidence applies to one pending receiver-side agent action. Course of action, workflow execution, and enforcement remain receiver-local.

\subsection{Trust, identity, and local policy}

KeyNote and macaroons support delegated authorization \cite{keynote99,birgisson2014macaroons}; TCX instead carries behavioral runtime evidence into a locally created decision. This avoids imported evidence becoming ambient enforcement authority, related to the confused-deputy concern \cite{hardy88}. Zero trust separates policy decision from enforcement \cite{nistZta}; SPIFFE federation authenticates foreign workloads \cite{spiffeFederation}; and RATS separates evidence, appraisal, verifier, and relying party \cite{rats}. TCOP applies those separations to cross-domain agent actions: it combines authenticated peer identity with receiver-local correlation and appraisal. Collaborative intrusion detection identifies relevant coordination and adversarial-participant questions \cite{bye2010}; the central audit here is only a bounded fact-equivalent control.

\subsection{Agent interoperability and tool protocols}

MCP standardizes client access to server tools and resources \cite{mcpAuthorization}; its gateway is the reference tool-call enforcement point here. A2A standardizes independent-agent discovery and collaboration \cite{a2aSpec}, and ANP addresses agent-network communication and discovery \cite{anpSpec}. They can supply transport or session context. Table~\ref{tab:protocol-comparison} summarizes the contract-level distinction: TCOP adds receipt-bound validation of signed runtime evidence against a pending receiver-side agent action, while leaving response selection and enforcement with the receiver.

\subsection{Agent and tool security}

Prompt injection and tool-integrated agents motivate the synthetic scenarios. Liu et al. benchmark prompt-injection attacks and defenses, while InjecAgent evaluates indirect injection against tool-integrated agents \cite{liu2024prompt,zhan2024injecagent}. TCOP does not compete with those detectors. It evaluates whether externally observed runtime evidence can affect a locally governed pending action in time.
